
Although establishing `structure-energy-property' relationships has been the overarching goal for materials modellings, most practices have been long focusing on performing calculations on a small sets of well-defined input atomistic models, to extract the corresponding energy and physical properties for each one of them at 0 K and generate some fundamental understandings therefrom. Throughout this review, based on \cref{eq:general_sampling_1} or \cref{eq:general_sampling_2}, it is one of our main intention to encourage the readers to establish a probabilistic or ensemble view and adopt the corresponding approach in the practice of materials modelling, as the probability theory forms the mathematical foundation for the development of modern ML/AI methodologies.

As we have discussed in previous sections, the probabilistic viewpoint is highly relevant to materials modelling. Fundamentally, realistic materials function at finite temperatures, and may be intrinsically disordered (such as alloys) or contains defects (such as oxygen vacancies) that are randomly distributed in the system. Performing PES sampling on material systems of interest helps to determine an ensemble of most likely accessible low-energy and stable structures under the thermodynamic conditions at which the materials will be applied. Hence computational modelling can be viewed as providing some probabilistic guidance on what the most likely structural states and corresponding physical properties are for a given material under certain thermodynamic conditions. In this regard, some of the physical properties and the theoretical approaches for computing them already are already consistent with this probabilistic view, including, for instances, (a) the configurational [\cref{eq:conf_entropy}] and vibrational [\cref{eq:vib_entropy}] entropies that are computed based on the Boltzmann distribution of microstates at given temperature $T$, (b) physical properties extracted from MD sampling [\cref{tab:properties_from_MD}], (c) stabilities and transition probabilities between different magnetic states [\cref{fig:mangetic_landscape}] and (d) the determinations of the strength of electron-phonon couplings and temperature-dependent renormalisation of electronic band energies, which needs to average over the thermally excited phononic and electronic states [\cref{sec:eph_coupling}]. At a more practical level, the numerical implementations of the theoretical methodologies can introduce noises and uncertainties, even though the underlying theories are formulated in a deterministic way, such as the random initialisations in gradient-based optimisers and the choices of different grid sizes when performing numerical integrations. These factors, sometimes, can lead to notable differences in the calculation outcomes, particually when performing statistical samplings on some very flat landscapes. 

Whilst most traditional materials modelling approaches were guided by strong physical priors or principles to determine the most likely observable(s) for a given structure, taking the above examples into considerations, a more probabilistics-based approach should be encouraged in future modelling practices, particularly to make it more compatible with the ML/AI-facilitated materials studies. This is because the latter takes an opposite approach to traditional materials modelling, as ML/AI techniques generally rely on large amounts of data to learn the relationships between the input structures and the output properties based on the probability distributions of both, and thus naturally takes into account the presences of noises in the data for making decisions with minimum or no human bias.  To facilitate the understanding of the probabilistic viewpoint in materials modelling and its connections to ML/AI, we provide a brief overview of the fundamentals of probability theory that are highly relevant to the development of modern ML/AI methodologies. Interested readers are referred to \citet{murphy2012machine} and \citet{goodfellow2016deep} for more detailed discussions.
